\subsection{Supporto alle dimensioni 9x9 e 19x19}

Un ulteriore obiettivo progettuale è stato supportare entrambe le dimensioni di scacchiera considerate: 9$\times$9 e 19$\times$19. A questo scopo, l'architettura è stata parametrizzata rispetto alla dimensione del goban, evitando di mantenere due implementazioni distinte.

La classe \texttt{GoTRMModel} utilizza tale parametro per determinare le dimensioni delle componenti dipendenti dal numero di intersezioni, in particolare la lunghezza della sequenza elaborata dal nucleo ricorsivo, la dimensione dell'embedding posizionale e il numero di azioni prodotte dalla Policy Head.

Le principali differenze tra le due configurazioni sono riassunte nella Tabella~\ref{tab:board_size_comparison}.

\begin{table}[H]
    \centering
    \begin{tabular}{l|c|c}
        \textbf{Caratteristica} & \textbf{9$\times$9} & \textbf{19$\times$19} \\
        \hline
        Intersezioni            & 81                  & 361                   \\
        Canali in ingresso      & 7                   & 7                     \\
        Token elaborati         & 81                  & 361                   \\
        Logits prodotti         & 82                  & 362                   \\
        Dimensione nascosta     & 256                 & 256
    \end{tabular}
    \caption{Confronto tra le configurazioni utilizzate per le due dimensioni di scacchiera.}
    \label{tab:board_size_comparison}
\end{table}

Sebbene la struttura architetturale sia condivisa, le due configurazioni utilizzano checkpoint distinti, ottenuti mediante processi di addestramento separati. Si tratta quindi di due modelli basati sulla stessa implementazione, ma configurati e addestrati specificamente per ciascuna dimensione di scacchiera.
